% Módulo 02a — Preliminares matemáticos e computacionais (R775+)
\chapter{Preliminares matemáticos e computacionais}
\section{Informação, entropia e exponenciação}
A teoria matemática da comunicação \cite{shannon1948} introduziu a entropia como medida de incerteza; a relação entre informação e suficiência foi formalizada por \cite{kullback1951}; e a conexão entre teoria da informação e mecânica estatística, com a distribuição exponencial como maximizadora de entropia sob média fixada, por \cite{jaynes1957}. É dessa linhagem que descende a ponderação exponencial normalizada: dados escores $f_i$, os pesos $w_i=e^{f_i}/\sum_j e^{f_j}$ formam a distribuição de máxima entropia compatível com a média dos escores. A interpretação probabilística de saídas de redes feedforward, com relações à recunohecimento estatístico de padrões, foi estabelecida por \cite{bridle1990}. Regras de escore estritamente próprias, que incentivam probabilidades calibradas, foram caracterizadas por \cite{gneiting2007}.
\section{Conjuntos convexos e ordens finitas}
Uma combinação $\sum w_ix_i$ com $w_i\geq0$ e $\sum w_i=1$ jamais escapa ao intervalo $[\min x_i,\max x_i]$: é o conteúdo de T2/T6, e a razão formal pela qual pesos devem somar um. Em conjuntos finitos não vazios sempre existe elemento maximal de qualquer função real (T3): finitude elimina a possibilidade de fuga ao infinito que invalida máximos em domínios infinitos. Listas ordenadas de forma decrescente têm a propriedade de que qualquer prefixo domina o sufixo (T9): é a anatomia do top-\textit{k}.
\section{Atenção neural e tradução automática}
A atenção produto-escalado com múltiplas cabeças foi introduzida para tradução automática \cite{vaswani2017}, com variantes efetivas exploradas em tradução neural \cite{luong2015}, pré-treino bidirecional \cite{devlin2019}, dependências longas \cite{dai2019}, benchmarks multitarefa \cite{wang2018}, compreensão de leitura em larga escala \cite{rajpurkar2016}, bases hierárquicas de imagens \cite{deng2009} e ferramental de processamento \cite{wolf2020}. Este trabalho herda apenas a forma --- ponderação normalizada de cabeças com decisão pelo máximo --- sem pesos treinados.
\section{Verificação, SMT e inspeção}
Bases axiomáticas para programação \cite{hoare1969} fundam a verificação dedutiva mecanizada em Lean \cite{moura2021} e Mathlib \cite{mathlib2020}; solucionadores SMT eficientes como o Z3 \cite{demoura2008} e computação simbólica como SymPy \cite{meurer2017} compõem o ecossistema de apoio. Inspeções de projeto e código reduzem erros antes da execução \cite{fagan1976}; complexidade ciclomática \cite{mccabe1976} e suítes de métricas orientadas a objetos \cite{chidamber1994} disciplinam a auditoria de código; processos racionais de projeto documentam decisões \cite{parnas1986}. Alocação de tarefas distribuídas segue desde o protocolo de rede de contratos \cite{smith1980}, e integração de conhecimento sob incerteza desde sistemas de quadro-negro \cite{erman1980}.
\section{Ponto flutuante e computação científica}
A aritmética de ponto flutuante exige atenção a arredondamento e casos especiais \cite{goldberg1991}; o manual de referência consolida algoritmos e análise \cite{muller2018}. A computação científica em Python apoia-se em arranjos \cite{harris2020} e algoritmos fundamentais \cite{virtanen2020}.
\section{Método empírico e estatística}
Experimentação disciplinada em engenharia de software \cite{basili1986,wohlin2012}, diretrizes de estudo de caso \cite{runeson2009} e relato sistemático \cite{page2021} estruturam o desenho; a declaração sobre valores-\textit{p} \cite{wasserstein2016}, a nova estatística por estimação \cite{cumming2014} e o bootstrap \cite{efron1979} fundamentam o relato por intervalos em vez de significância pontual. Reprodutibilidade como problema de campo \cite{osc2015,collberg2016,ioannidis2005}, dados documentados \cite{gebru2021}, modelos documentados \cite{mitchell2019} e dados geridos por FAIR \cite{wilkinson2016} justificam a trilha publicada.
\section{Cálculo dito passo a passo}
\textit{Normalização.} Com $S=\sum_j e^{f_j}>0$, cada quociente $e^{f_i}/S$ é positivo e a soma telescopa em $S/S=1$. \textit{Deslocamento.} Subtrair $c$ multiplica numerador e denominador por $1/e^c$, que cancela: a decisão independe do nível. \textit{Temperatura.} Para $\tau_1<\tau_2$ e maximizador único $i^*$, a razão de pesos cresce porque cada termo concorrente $e^{(f_j-f_{i^*})/\tau}$ decai mais rápido no $\tau$ menor --- a Tabela~\ref{tab:temp} mostra $0{,}5065\to0{,}9817$. \textit{Overflow.} Deslocado pelo máximo, todo expoente é $\leq0$ e cada exponencial fica em $(0,1]$: nenhum estouro possível.
\section{Top-\textit{k} e complexidade}
Extrair os $k$ maiores por heap custa $O(n\log k)$ contra $O(n\log n)$ da ordenação total; o desempate por identificador torna a decisão determinística. A garantia independe do algoritmo: sob ordem decrescente, o prefixo domina o sufixo.
